% AUTO-GENERATED BY Experiment pipeline. DO NOT EDIT BY HAND.
\section*{Bộ Plot Mới Từ Experiment (Nguồn Tự Viết)}

\subsection*{EXP Hình 1: Hiệu ứng gộp}
\begin{figure}[H]
  \centering
  \includegraphics[width=\textwidth]{EXP_Figure1_Pooled_Effects.png}
  \caption{Ước lượng gộp với khoảng tin cậy và dị hợp cho ba đối sánh.}
\end{figure}
\noindent\textbf{Insight:} Cộng hưởng mạnh âm ($g=-0.448$), Tăng cường cho con người dương ($g=0.563$), và Tăng cường cho AI không có ý nghĩa ở mức gộp ($p=0.780$).

\subsection*{EXP Hình 2: Sai lệch công bố}
\begin{figure}[H]
  \centering
  \includegraphics[width=\textwidth]{EXP_Figure2_Bias_Diagnostics.png}
  \caption{Mức bằng chứng sai lệch công bố (Egger/Begg) và số nghiên cứu thiếu (trim-and-fill).}
\end{figure}
\noindent\textbf{Insight:} Cả cường độ sai lệch và quy mô thiếu hụt đều đáng kể ở ba đối sánh.

\subsection*{EXP Hình 3: Biên độ điều tiết}
\begin{figure}[H]
  \centering
  \includegraphics[width=\textwidth]{EXP_Figure3_Subgroup_Range.png}
  \caption{Biên độ dao động của hiệu ứng giữa mức cao nhất và thấp nhất theo điều tiết.}
\end{figure}
\noindent\textbf{Insight:} Biên độ lớn hàm ý tính phụ thuộc bối cảnh mạnh và nguy cơ đảo dấu hiệu ứng.

\subsection*{EXP Hình 4: Hồ sơ ngành}
\begin{figure}[H]
  \centering
  \includegraphics[width=\textwidth]{EXP_Figure4_Industry_Profile.png}
  \caption{Hiệu ứng theo ngành cho ba đối sánh chính.}
\end{figure}
\noindent\textbf{Insight:} Riêng Tăng cường cho AI dao động từ Kinh doanh ($g=-0.732$) đến Truyền thông ($g=1.011$).

\subsection*{EXP Hình 5: Phân bố và ngoại lệ}
\begin{figure}[H]
  \centering
  \includegraphics[width=\textwidth]{EXP_Figure5_Distribution_Quality.png}
  \caption{Chẩn đoán phân bố và ngoại lệ, nhấn mạnh đuôi phải của AI Augmentation.}
\end{figure}
\subsection*{EXP Hình 6: Mất cân bằng moderator}
\begin{figure}[H]
  \centering
  \includegraphics[width=\textwidth]{EXP_Figure6_Moderator_Imbalance.png}
  \caption{Mức mất cân bằng nhóm phổ biến nhất theo moderator.}
\end{figure}

\subsection*{EXP Hình 7: Xu hướng theo năm}
\begin{figure}[H]
  \centering
  \includegraphics[width=\textwidth]{EXP_Figure7_Temporal_Trend.png}
  \caption{Xu hướng effect size mô tả theo năm (kèm vùng bất định 95\% xấp xỉ).}
\end{figure}
\noindent\textbf{Insight:} AI Augmentation có biên độ dao động mean theo năm $\Delta g=1.118$ (mô tả), nhấn mạnh nguy cơ drift theo giai đoạn.

\subsection*{EXP Hình 8: Cấu trúc dấu effect size}
\begin{figure}[H]
  \centering
  \includegraphics[width=\textwidth]{EXP_Figure8_Sign_Profile.png}
  \caption{Tỷ trọng effect size theo ba vùng: âm rõ rệt, trung tính, dương rõ rệt.}
\end{figure}
\noindent\textbf{Insight:} Ở AI Augmentation, tỷ trọng âm rõ rệt=49.3\%, trung tính=9.0\%, dương rõ rệt=41.7\%.

\subsection*{EXP Hình 9: Mô tả vs pooled}
\begin{figure}[H]
  \centering
  \includegraphics[width=\textwidth]{EXP_Figure9_Descriptive_vs_Pooled.png}
  \caption{Đối chiếu mean/median mô tả với ước lượng pooled từ mô hình 3 mức.}
\end{figure}
\noindent\textbf{Insight:} Với AI Augmentation, mean=0.197, median=-0.141, pooled=0.058, hàm ý mô tả thuần túy có thể bị kéo lệch bởi đuôi phân bố.

\noindent\textit{Lưu ý:} Toàn bộ hình ở section này được sinh trực tiếp từ mã nguồn trong thư mục Experiment.
